% =====================================================================
%  Couverture, page de titre et plat verso du dictionnaire de poche.
%
%  Le modele est la maquette fournie par l'editeur de cette edition
%  (Dicionario.pages, A4) : plat recto bleu a lettrage noir et etoile
%  blanche, page de titre blanche a lettrage noir et etoile bleue, plat
%  verso bleu uni. Tout ce qui suit — corps des caracteres, hauteurs de
%  ligne, diametre de l'embleme — est RELEVE sur ce modele au 300e de
%  pouce, puis rapporte a la page de poche. Les valeurs mesurees sont
%  donnees en regard, pour qu'on puisse refaire le calcul.
% =====================================================================

\usepackage{tikz}
\usepackage{eso-pic}

% Le bleu est l'AZUR, #007FFF — le bleu de reference, a mi-chemin exact du
% bleu et du cyan sur la roue sRGB (teinte 210 degres, saturation e valeur
% pleines). C'est celui des deux plats, du disque et du lettrage ID de la
% page de titre : une seule couleur pour tout le livre.
%
% Il ne vient pas d'un releve sur la maquette : celle-ci est ecrite dans
% l'espace Display P3, ou son bleu vaut (0,2549 ; 0,5728 ; 0,9688), soit
% #0094FF une fois converti en sRGB d'apres le profil enferme dans le
% fichier. L'azur en est proche — vingt et un points de vert plus bas — et
% c'est une couleur NOMMEE, non le residu d'une conversion.
\definecolor{idoblua}{RGB}{0,127,255}

% Futura n'est pas libre ; Jost* en est la reprise sous licence ouverte.
% Chargee PAR CHEMIN, non par nom : le document se compose a l'identique
% sur une machine ou la police n'est pas installee. Voir polices/LISEZ-MOI.md.
\newfontfamily\kovrilfonto{Jost}[Path=polices/, Extension=.ttf,
  UprightFont=*-Medium, BoldFont=*-Bold, ItalicFont=*-MediumItalic]

% --- Corps des caracteres --------------------------------------------
% Le modele est compose en Futura sur une page A4 (595,28 pt) : 62 pt pour
% DICIONARIO, 38 pt pour le sous-titre, 204 pt pour IDO, 29 pt pour la
% signature. Deux corrections les rapportent a la page de poche :
%
%   1. le rapport des largeurs de page, 110 mm / 210 mm ;
%   2. le rapport des hauteurs de capitale, Futura 0,754 em contre
%      Jost 0,700 em, soit 1,0770 — sans quoi le lettrage, sur une
%      couverture qui n'est presque que capitales, paraitrait plus petit
%      que le modele a corps egal.
%
% Les deux se composent en un seul facteur, applique a \paperwidth.
\newcommand{\korpDICIONARIO}{0.11217\paperwidth}   % 62 pt sur A4
\newcommand{\korpSUBTITOLO} {0.06875\paperwidth}   % 38 pt
\newcommand{\korpIDO}       {0.36908\paperwidth}   % 204 pt
\newcommand{\korpNOMO}      {0.05247\paperwidth}   % 29 pt

% --- Hauteurs de ligne ------------------------------------------------
% Position du HAUT de l'encre de chaque ligne, comptee depuis le bord
% superieur de la feuille. Trois decisions les fixent.
%
% 1. L'ECHELLE. Les distances verticales suivent le rapport des LARGEURS de
%    page, 110/210 = 0,5238, et non celui des corps. Les deux different : le
%    corps a ete grossi de 1,0770 pour egaliser la hauteur de capitale de
%    Jost* et celle de la Futura, et c'est la capitale, non le cadratin, que
%    l'oeil mesure. Une distance verticale mise a l'echelle des capitales est
%    donc mise a l'echelle des largeurs — les deux facteurs se compensent
%    exactement, et le diametre de l'embleme le verifie : 204 x 0,5641 x 0,700
%    = 153,8 x 0,5238.
%
% 2. L'INTERLIGNAGE DU SOUS-TITRE. Le modele porte 50,1 pt entre la premiere
%    et la deuxieme ligne, 56,2 entre la deuxieme et la troisieme, la ou les
%    trois sont du meme corps. C'est un reglage de paragraphe reste inegal,
%    non une intention : on prend la moyenne, 53,15 pt, et « di la » descend
%    de 3,05 pt sur le gabarit A4, soit 1,6 pt sur la page de poche.
%
% 3. LA POSE DU BLOC. Les deux marges qui restent sont partagees en SECTION
%    DOREE, la petite en haut : 81,75 pt contre 132,27, dont le rapport est
%    1/phi = 0,6180. Le modele les partageait 87,8 contre 188,6, soit 0,466 —
%    proportion qui allait a une page A4 mais qui, reportee sur une page plus
%    haute que large, poussait la composition trop haut.
%
%    Le controle vient de l'embleme, qui est ce que l'oeil pese : sa section
%    doree le centre a 0,5297 de la hauteur de page, la ou le modele le place
%    a 0,5323 — deux millimes d'ecart. L'ancienne pose le mettait a 0,5034,
%    c'est-a-dire au milieu mort de la page.
\newcommand{\posTITOLO}  {0.160218\paperheight}
\newcommand{\posRADIKI}  {0.235673\paperheight}
\newcommand{\posDILA}    {0.290237\paperheight}
\newcommand{\posLINGUO}  {0.344801\paperheight}
\newcommand{\posIDO}     {0.445613\paperheight}
\newcommand{\posNOMO}    {0.675162\paperheight}
\newcommand{\posAKADEMIO}{0.715404\paperheight}

% --- L'embleme -------------------------------------------------------
% Un disque et une etoile a six branches. Elle est REGULIERE, et se
% construit :
%
%   1. un triangle equilateral au centre, une pointe en bas — ses trois
%      sommets sont les trois petites branches ;
%   2. sur chacun de ses cotes, une longue pointe dressee vers l'exterieur,
%      dont le sommet touche le cercle ;
%   3. la longue pointe du haut a la meme hauteur que le triangle ;
%   4. chaque longue pointe a pour base le TIERS du cote qui la porte.
%
% La regle 4 en recouvre une autre, plus parlante, et probablement celle du
% symbole d'origine : le bord d'une longue pointe, PROLONGE vers l'interieur,
% passe par le milieu d'un des deux autres cotes du triangle. Les deux regles
% donnent la meme figure, et c'est encore t = 1/2 qui les accorde — la
% construction par les milieux donne a la base (2-t)/(4+t) du cote, soit un
% tiers exactement lorsque t vaut un demi.
%
% La condition 3 suffit a tout fixer. En prenant le rayon du disque pour
% unite et t pour rayon circonscrit du triangle, la hauteur du triangle vaut
% 3t/2 et celle d'une pointe 1 moins le rayon inscrit, soit 1 - t/2 ; les
% egaler donne t = 1/2. Le triangle tient donc dans le demi-disque, son cote
% vaut V3/2, son rayon inscrit 1/4, et le tiers de son cote V3/6 — d'ou tous
% les nombres ci-dessous, qui sont des multiples de V3/12 et de 1/4.
%
% La figure est invariante par rotation d'un tiers de tour : les trois
% grandes pointes sont sur le cercle, les trois petites a mi-rayon.
%
% Confrontee a la maquette, dont l'etoile est dessinee a la main, la figure
% construite s'en ecarte de 2,4 % du rayon au plus — l'etoile de la maquette
% est posee 1,4 % trop bas dans son disque, et ses longues pointes ont a la
% base 0,30 fois le cote et non un tiers.
%
% Coordonnees en unites de RAYON, origine au centre du disque, y vers le
% haut. La ligne de base de la figure passe par le centre du disque.
\newlength{\stelrayo}
\newcommand{\idoemblemo}[3]{% #1 diametre, #2 couleur du disque, #3 de l'etoile
  \setlength{\stelrayo}{0.5\dimexpr#1\relax}%
  \begin{tikzpicture}[x=\stelrayo, y=\stelrayo, baseline={(0,0)}]
    \fill[#2] (0,0) circle[radius=1];
    \fill[#3]
         (            0 ,  1   )   % longue pointe du haut
      -- ({-sqrt(3)/12 },  0.25)   % base de la pointe du haut
      -- ({-sqrt(3)/4  },  0.25)   % sommet gauche du triangle
      -- ({-sqrt(3)/6  },  0   )   % base de la pointe bas-gauche
      -- ({-sqrt(3)/2  }, -0.5 )   % longue pointe bas-gauche
      -- ({-sqrt(3)/12 }, -0.25)   % base de la pointe bas-gauche
      -- (            0 , -0.5 )   % sommet bas du triangle
      -- ({ sqrt(3)/12 }, -0.25)
      -- ({ sqrt(3)/2  }, -0.5 )   % longue pointe bas-droite
      -- ({ sqrt(3)/6  },  0   )
      -- ({ sqrt(3)/4  },  0.25)   % sommet droit du triangle
      -- ({ sqrt(3)/12 },  0.25)
      -- cycle;
  \end{tikzpicture}}

% Le mot IDO : les deux premieres lettres sont composees, la troisieme est
% l'embleme. Sur le modele le disque vaut 1,0651 fois la hauteur de
% capitale et se centre sur elle, et il suit le D a 1,4 pt.
%
% Les lettres prennent la couleur du DISQUE et non celle du texte : sur le
% plat elles sont blanches quand le titre est noir, sur la page de titre
% elles sont bleues. Le mot fait bloc avec l'embleme, il ne fait pas bloc
% avec le lettrage.
\newlength{\kapalto}
% Le diametre, mesure une fois pour toutes a l'ouverture du document : les
% deux plats et la page de titre portent le MEME embleme, et il ne doit pas
% dependre de l'endroit ou on le pose.
\newlength{\emblemdiam}
\AtBeginDocument{{\kovrilfonto\bfseries\fontsize{\korpIDO}{\korpIDO}\selectfont
  \settoheight{\kapalto}{ID}\global\emblemdiam=1.0651\kapalto}}
\newcommand{\idogrupo}[2]{% #1 couleur du disque, #2 couleur de l'etoile
  {\kovrilfonto\bfseries\fontsize{\korpIDO}{\korpIDO}\selectfont
   \settoheight{\kapalto}{ID}%
   \textcolor{#1}{ID}\kern0.0045\paperwidth
   \raisebox{0.5\kapalto}{\idoemblemo{\emblemdiam}{#1}{#2}}}}

% --- Pose des lignes --------------------------------------------------
% \supre{hauteur}{contenu} : le contenu centre sur la largeur de la
% feuille, le HAUT de son encre a la hauteur dite. \raisebox{-\height}
% descend la boite de sa propre hauteur, ce qui amene son sommet sur la
% ligne de base du \put.
\newcommand{\supre}[2]{%
  \put(0,-#1){\makebox[\paperwidth][c]{\raisebox{-\height}{#2}}}}

% Le lettrage, identique sur les deux pages ; seules changent les couleurs.
\newcommand{\kovrilteksto}[3]{% #1 encre, #2 disque, #3 etoile
  \AddToShipoutPictureFG*{\AtPageUpperLeft{%
    \setlength{\unitlength}{1pt}\color{#1}\kovrilfonto
    \supre{\posTITOLO}{\bfseries\fontsize{\korpDICIONARIO}{\korpDICIONARIO}
      \selectfont DICIONARIO}%
    \supre{\posRADIKI}{\bfseries\fontsize{\korpSUBTITOLO}{\korpSUBTITOLO}
      \selectfont de la 10.000 RADIKI}%
    \supre{\posDILA}{\bfseries\fontsize{\korpSUBTITOLO}{\korpSUBTITOLO}
      \selectfont di la}%
    \supre{\posLINGUO}{\bfseries\fontsize{\korpSUBTITOLO}{\korpSUBTITOLO}
      \selectfont linguo universala}%
    \supre{\posIDO}{\idogrupo{#2}{#3}}%
    \supre{\posNOMO}{\fontsize{\korpNOMO}{\korpNOMO}\selectfont
      Marcelo Persiko \textit{(Marcel Pesch)}}%
    \supre{\posAKADEMIO}{\fontsize{\korpNOMO}{\korpNOMO}\selectfont
      ex-membro di la Akademio}%
  }}}

\newcommand{\fondo}[1]{%
  \AddToShipoutPictureBG*{\AtPageLowerLeft{%
    \color{#1}\rule{\paperwidth}{\paperheight}}}}

% --- Les trois pages --------------------------------------------------
% Le plat recto : fond bleu, lettrage noir, disque blanc, etoile bleue.
\newcommand{\platrekto}{%
  \fondo{idoblua}\kovrilteksto{black}{white}{idoblua}%
  \thispagestyle{empty}\null\clearpage}

% La page de titre : fond blanc, lettrage noir, disque bleu, etoile blanche.
\newcommand{\titolpagino}{%
  \kovrilteksto{black}{idoblua}{white}%
  \thispagestyle{empty}\null\clearpage}

% Le plat verso : l'embleme seul sur le bleu. Meme diametre qu'au plat recto
% — c'est la meme marque, et la repeter a la meme taille fait rimer les deux
% plats. Le disque est blanc et l'etoile bleue, comme au recto : sur un fond
% bleu, c'est le disque qui doit porter la lumiere.
%
% La hauteur n'est pas le milieu geometrique : un signe seul y parait BAS,
% et il faut le centre optique. Celui que l'oeil demande ici se trouve etre
% le MIROIR du centre de l'embleme du plat recto autour de l'axe horizontal
% de la page — 0,5297 au recto, 1 moins 0,5297 au verso. Les deux plats
% tiennent donc la marque a des hauteurs symetriques : un peu au-dessous du
% milieu quand une composition l'accompagne, un peu au-dessus quand elle est
% seule.
\newcommand{\posVERSO}{0.470298\paperheight}
\newlength{\posverso}
\newcommand{\platverso}{%
  \fondo{idoblua}%
  \AddToShipoutPictureFG*{\AtPageUpperLeft{%
    \setlength{\unitlength}{1pt}%
    \setlength{\posverso}{\posVERSO}%
    \addtolength{\posverso}{-0.5\emblemdiam}%
    \put(0,-\posverso){\makebox[\paperwidth][c]{\raisebox{-\height}{%
      \idoemblemo{\emblemdiam}{white}{idoblua}}}}}}%
  \thispagestyle{empty}\null\clearpage}
